% =====================================================================
% ANHANGTEIL -- eigene Nummerierung (Anhang 3, Anhang 5),
% Bezeichnungen analog zu den zugehoerigen Abschnitten
% =====================================================================

% ---------------------------------------------------------------------
\setzeTeil{Anhang 3}
\section{Tarifliche monatliche Beitragsraten}\label{sec:anhang3}

Tarifbeiträge gemäß den Gleichungen~\eqref{eq:netto},
\eqref{eq:brutto} und \eqref{eq:bkinder} (Einzelversicherung, ohne
gesetzlichen Beitragszuschlag und ohne Risikozuschläge; Stand:
Änderungsnr.~17, BAP 01/2025; Endalter der Tafel: 110).
Höchsteintrittsalter des (geschlossenen) Tarifs: 75~Jahre; Werte
darüber dienen tariftechnischen Zwecken.

\begin{center}
\begin{tabular}{lcc}
\toprule
Beobachtungseinheit & $P$ (EUR p.\,a.) & $b$ (EUR/Monat) \\
\midrule
Kinder (0--15) & \KKinder & \BKinder \\
männliche Jugendliche (16--20) & \KJugendM & \BJugendM \\
weibliche Jugendliche (16--20) & \KJugendF & \BJugendF \\
\bottomrule
\end{tabular}
\end{center}

% Automatisch erzeugt von kalkulation/pipeline.py
\begin{small}
\begin{longtable}{rrrrr}
\caption{Tarifliche monatliche Beitragsraten $b_x$ (EUR/Monat, Einzelversicherung) und Jahresnettobeiträge $P_x$ (EUR p.\,a.)}\label{tab:beitraege}\\
\toprule
$x$ & $P_x$ M & $b_x$ M & $P_x$ F & $b_x$ F \\
\midrule
\endfirsthead
\caption[]{Tarifliche monatliche Beitragsraten $b_x$ (EUR/Monat, Einzelversicherung) und Jahresnettobeiträge $P_x$ (EUR p.\,a.) (Fortsetzung)}\\
\toprule
$x$ & $P_x$ M & $b_x$ M & $P_x$ F & $b_x$ F \\
\midrule
\endhead
\bottomrule
\endfoot
21 & 2446,25 & 334,14 & 3353,01 & 431,28 \\
22 & 2499,95 & 338,71 & 3440,18 & 438,77 \\
23 & 2560,47 & 343,75 & 3525,95 & 445,96 \\
24 & 2627,91 & 349,26 & 3608,77 & 452,70 \\
25 & 2701,85 & 355,41 & 3687,31 & 459,13 \\
26 & 2782,99 & 362,25 & 3761,89 & 465,25 \\
27 & 2870,77 & 369,74 & 3833,04 & 471,13 \\
28 & 2958,79 & 377,30 & 3902,55 & 476,93 \\
29 & 3048,02 & 385,01 & 3971,98 & 482,77 \\
30 & 3138,27 & 392,85 & 4041,77 & 488,70 \\
31 & 3230,42 & 400,90 & 4112,97 & 494,81 \\
32 & 3323,97 & 409,12 & 4185,40 & 501,08 \\
33 & 3420,66 & 417,66 & 4259,28 & 507,53 \\
34 & 3520,41 & 426,53 & 4336,84 & 514,35 \\
35 & 3624,94 & 435,83 & 4419,04 & 521,61 \\
36 & 3734,18 & 445,56 & 4504,90 & 529,20 \\
37 & 3848,64 & 455,76 & 4593,95 & 537,09 \\
38 & 3968,57 & 466,47 & 4685,77 & 545,24 \\
39 & 4095,91 & 477,86 & 4781,68 & 553,76 \\
40 & 4229,30 & 489,79 & 4880,69 & 562,57 \\
41 & 4369,33 & 502,32 & 4982,85 & 571,66 \\
42 & 4515,48 & 515,40 & 5087,93 & 581,02 \\
43 & 4667,48 & 529,01 & 5195,91 & 590,64 \\
44 & 4824,83 & 543,10 & 5306,52 & 600,50 \\
45 & 4987,04 & 557,63 & 5419,47 & 610,56 \\
46 & 5153,63 & 572,54 & 5534,07 & 620,77 \\
47 & 5324,18 & 587,80 & 5649,75 & 631,07 \\
48 & 5498,32 & 603,37 & 5765,93 & 641,40 \\
49 & 5675,71 & 619,23 & 5882,09 & 651,73 \\
50 & 5855,92 & 653,33 & 5997,90 & 662,02 \\
51 & 6038,49 & 651,63 & 6113,13 & 672,25 \\
52 & 6222,96 & 668,10 & 6227,66 & 682,41 \\
53 & 6408,76 & 684,67 & 6341,53 & 692,51 \\
54 & 6595,19 & 701,28 & 6454,86 & 702,55 \\
55 & 6781,56 & 717,94 & 6567,84 & 712,61 \\
56 & 6968,43 & 734,72 & 6681,55 & 722,79 \\
57 & 7156,37 & 751,67 & 6796,94 & 733,19 \\
58 & 7346,01 & 768,86 & 6914,80 & 743,89 \\
59 & 7538,04 & 786,36 & 7035,80 & 754,94 \\
60 & 7733,21 & 804,17 & 7160,48 & 766,34 \\
61 & 7930,94 & 822,24 & 7288,51 & 778,06 \\
62 & 8130,59 & 840,52 & 7419,40 & 790,06 \\
63 & 8331,39 & 858,94 & 7552,55 & 802,30 \\
64 & 8532,49 & 877,43 & 7687,33 & 814,70 \\
65 & 8733,15 & 883,84 & 7823,29 & 813,43 \\
66 & 8932,67 & 902,27 & 7960,03 & 826,07 \\
67 & 9130,66 & 920,62 & 8097,28 & 838,79 \\
68 & 9326,82 & 938,87 & 8234,62 & 851,57 \\
69 & 9520,91 & 957,01 & 8371,46 & 864,36 \\
70 & 9712,32 & 974,98 & 8507,02 & 877,08 \\
71 & 9900,12 & 992,72 & 8640,38 & 889,67 \\
72 & 10083,32 & 1010,14 & 8770,36 & 902,02 \\
73 & 10260,83 & 1027,15 & 8895,77 & 914,04 \\
74 & 10431,66 & 1043,68 & 9015,39 & 925,62 \\
75 & 10594,87 & 1059,65 & 9128,22 & 936,69 \\
76 & 10749,73 & 1068,94 & 9233,60 & 942,66 \\
77 & 10893,24 & 1076,69 & 9331,01 & 947,54 \\
78 & 11025,77 & 1082,86 & 9420,02 & 951,25 \\
79 & 11146,52 & 1087,33 & 9500,42 & 953,72 \\
80 & 11254,95 & 1089,96 & 9572,25 & 954,89 \\
81 & 11350,68 & 1098,42 & 9633,89 & 960,33 \\
82 & 11431,96 & 1105,59 & 9686,23 & 964,95 \\
83 & 11497,63 & 1111,39 & 9729,26 & 968,75 \\
84 & 11547,19 & 1115,76 & 9763,41 & 971,76 \\
85 & 11581,41 & 1118,78 & 9789,37 & 974,05 \\
86 & 11601,82 & 1120,59 & 9808,09 & 975,71 \\
87 & 11611,36 & 1121,43 & 9820,91 & 976,84 \\
88 & 11613,71 & 1121,63 & 9828,80 & 977,53 \\
89 & 11613,81 & 1121,64 & 9832,75 & 977,88 \\
90 & 11613,84 & 1121,65 & 9833,88 & 977,98 \\
91 & 11613,81 & 1121,64 & 9833,77 & 977,97 \\
92 & 11613,79 & 1121,64 & 9833,70 & 977,97 \\
93 & 11613,65 & 1121,63 & 9833,74 & 977,97 \\
94 & 11613,69 & 1121,63 & 9833,76 & 977,97 \\
95 & 11613,77 & 1121,64 & 9833,85 & 977,98 \\
96 & 11613,59 & 1121,62 & 9833,81 & 977,98 \\
97 & 11613,39 & 1121,61 & 9833,95 & 977,99 \\
98 & 11613,50 & 1121,62 & 9833,86 & 977,98 \\
99 & 11613,69 & 1121,63 & 9833,96 & 977,99 \\
100 & 11613,64 & 1121,63 & 9833,81 & 977,98 \\
101 & 11613,53 & 1121,62 & 9833,73 & 977,97 \\
102 & 11613,40 & 1121,61 & 9833,76 & 977,97 \\
103 & 11613,41 & 1121,61 & 9833,76 & 977,97 \\
104 & 11613,71 & 1121,64 & 9834,01 & 977,99 \\
105 & 11613,35 & 1121,60 & 9833,77 & 977,97 \\
106 & 11613,49 & 1121,62 & 9833,85 & 977,98 \\
107 & 11613,88 & 1121,65 & 9833,86 & 977,98 \\
108 & 11613,92 & 1121,65 & 9834,28 & 978,02 \\
109 & 11614,00 & 1121,66 & 9834,15 & 978,01 \\
110 & 11613,63 & 1121,63 & 9833,92 & 977,99 \\
\end{longtable}
\end{small}


% ---------------------------------------------------------------------
\setzeTeil{Anhang 5}
\setcounter{section}{4}
\section{Beispielrechnung}\label{sec:anhang5}

Die folgende Beispielrechnung vollzieht die Beitragsberechnung für
eine \textbf{männliche Person mit Eintrittsalter $x_e = \BspAlter$}
Schritt für Schritt nach. Alle Werte folgen aus den Anlagen dieser TB
und den Rundungskonventionen des Abschnitts~\ref{sec:rundung}.

\bigskip\noindent\textbf{Schritt 1: Rechnungsgrundlagen (Anlagen 2.1, 2.2 und 2.3.3)}
\begin{align*}
i &= \ZinsM\,\% \quad\Rightarrow\quad v = \tfrac{1}{1{,}027} = 0{,}973710\dots\\
k_{\BspAlter} &= \Bspkx, \qquad
K_{\BspAlter} = k_{\BspAlter}\cdot G^{M} = \Bspkx \cdot \GMann = \BspKx\ \text{EUR}\\
q_{\BspAlter} &= \Bspqx, \qquad w_{\BspAlter} = \Bspwx\\
DQ_{\BspAlter} &= v\cdot\bigl(1 - q_{\BspAlter} - w_{\BspAlter}\bigr) = \BspDQx
\end{align*}

\bigskip\noindent\textbf{Schritt 2: Barwerte (Anlage 3.1, Rekursionen
\eqref{eq:rekA}/\eqref{eq:reka})}
\begin{align*}
A_{\BspAlter} &= k_{\BspAlter} + DQ_{\BspAlter}\cdot A_{44}
= \Bspkx + \BspDQx\cdot\BspAxE = \BspAx\\
a_{\BspAlter} &= 1 + DQ_{\BspAlter}\cdot a_{44}
= 1 + \BspDQx\cdot\BspaxE = \Bspax
\end{align*}

\bigskip\noindent\textbf{Schritt 3: Jahresnettobeitrag (Gleichung \eqref{eq:netto})}
\begin{equation*}
P_{\BspAlter} = \frac{G^{M}\cdot A_{\BspAlter}}{a_{\BspAlter}}
= \frac{\GMann\cdot\BspAx}{\Bspax} = 4676{,}48\ \text{EUR p.\,a.}
\end{equation*}

\bigskip\noindent\textbf{Schritt 4: Monatlicher Bruttobeitrag (Gleichung \eqref{eq:brutto})}

\noindent Mit $\gamma_{\BspAlter} = \BspGamma$~EUR, $\sigma + \Omega = 5{,}7\,\%$
und $\alpha_{\BspAlter} = \AlphaZ$:
\begin{equation*}
b_{\BspAlter}
= \frac{\GMann\cdot\BspAx + \BspGamma\cdot\Bspax}
       {12\cdot(1-0{,}057)\cdot\Bspax - 3{,}5}
= \Bspbx\ \text{EUR/Monat}.
\end{equation*}
Der monatliche Zahlbetrag bis zum Alter 60 beträgt einschließlich des
gesetzlichen Beitragszuschlags $1{,}1\cdot b_{\BspAlter}$
(Gleichung~\eqref{eq:zahlbeitrag}).

\bigskip\noindent\textbf{Schritt 5: Alterungsrückstellung (Gleichung \eqref{eq:ar})}
\begin{center}
\begin{tabular}{lr}
\toprule
Zeitpunkt & $V_{\BspAlter+m}(b_{\BspAlter})$ (EUR) \\
\midrule
Vertragsbeginn ($m=0$; $\approx -\alpha_{\BspAlter}\cdot b_{\BspAlter}$) & \BspVnull \\
nach 5 Jahren & \BspVfuenf \\
nach 10 Jahren & \BspVzehn \\
nach 20 Jahren & \BspVzwanzig \\
\bottomrule
\end{tabular}
\end{center}
Der Anfangswert folgt aus Gleichung~\eqref{eq:ar}; die Identität
$V_{x_e} = -\alpha_{x_e}\cdot b_{x_e}$ gilt wegen der kaufmännischen
Rundung des Beitrags auf 0{,}01~EUR nur näherungsweise
(Abschnitt~\ref{sec:rundung}). Die Rückstellung ist ab dem
Versicherungsjahr $m = \BspNegM$ nichtnegativ; der Rückgang beim
Übergang vom Alter 64 zum Alter 65 ($m = 21 \to 22$) beruht auf dem
Stufenwechsel der Zuschläge (Festlegung in
Abschnitt~\ref{sec:tarifar}). Der vollständige Verlauf der ersten
30 Jahre (Endalter der Rückstellungsbewertung: 110):

% Automatisch erzeugt von kalkulation/pipeline.py
\begin{small}
\begin{longtable}{rrr}
\caption{Beispiel: Verlauf der gezillmerten Alterungsrückstellung, Männer, Eintrittsalter 43 (EUR)}\label{tab:bspar}\\
\toprule
$m$ & erreichtes Alter & $V_{43+m}(b_{43})$ \\
\midrule
\endfirsthead
\caption[]{Beispiel: Verlauf der gezillmerten Alterungsrückstellung, Männer, Eintrittsalter 43 (EUR) (Fortsetzung)}\\
\toprule
$m$ & erreichtes Alter & $V_{43+m}(b_{43})$ \\
\midrule
\endhead
\bottomrule
\endfoot
0 & 43 & -1852,15 \\
1 & 44 & 984,98 \\
2 & 45 & 3944,67 \\
3 & 46 & 7023,61 \\
4 & 47 & 10220,05 \\
5 & 48 & 13534,18 \\
6 & 49 & 16967,29 \\
7 & 50 & 20520,24 \\
8 & 51 & 24194,03 \\
9 & 52 & 27990,26 \\
10 & 53 & 31909,47 \\
11 & 54 & 35950,32 \\
12 & 55 & 39848,70 \\
13 & 56 & 43538,28 \\
14 & 57 & 46952,71 \\
15 & 58 & 50028,06 \\
16 & 59 & 52704,08 \\
17 & 60 & 55268,65 \\
18 & 61 & 57697,47 \\
19 & 62 & 59968,22 \\
20 & 63 & 62058,75 \\
21 & 64 & 63944,07 \\
22 & 65 & 63454,09 \\
23 & 66 & 64939,96 \\
24 & 67 & 66181,92 \\
25 & 68 & 67172,70 \\
26 & 69 & 67906,30 \\
27 & 70 & 68384,25 \\
28 & 71 & 68596,94 \\
29 & 72 & 68537,22 \\
30 & 73 & 68199,01 \\
\end{longtable}
\end{small}

